% Abstract and keywords for the IEEE TC Graph-Centric Computing special issue.

% Title: OrbitGraph: Graph-Centric Software Defined Routing for LEO Satellite Constellations

\begin{abstract}
Low Earth orbit constellations form \emph{dynamic graphs}: links and weights change
as satellites move and ground-satellite contacts hand over.
Distributed OSPF must reconverge after each topology change, which can interrupt
end-to-end forwarding for multiple seconds.
We propose OrbitGraph, a graph-centric software-defined routing scheme that
recomputes shortest paths centrally and pushes only changed kernel routes
\emph{before} the old ground-satellite link is torn down.
Using container-based emulation from 27 to 100 satellites (including 2 ground stations), 
we measure control-plane cost and data-plane availability with a continuous outage 
probe aligned to routing events.
On coverage-preserving handovers, OrbitGraph yields 0\,ms data-plane outage and
0\% packet loss, while OSPF incurs multi-second black holes (up to ${\sim}15$\,s).
OrbitGraph pushes 524 kernel routes at the largest constellation versus 10\,504
at full initialization, and Dijkstra recomputation remains in the millisecond range.
When the constellation is too sparse to keep a ground station in view, a coverage
gap severs the ground-to-ground path for both protocols, bounding where
control-plane update ordering helps.
\end{abstract}

\begin{IEEEkeywords}
Dynamic graphs,
LEO satellite networks,
software-defined networking,
topology handover,
incremental routing,
network emulation,
graph-centric systems,
data-plane availability,
routing scalability
\end{IEEEkeywords}
